% «Más ejemplos resueltos» · Comprensión lectora (libros 3 y 4): un fragmento corto, una pregunta y su respuesta razonada.
\newcommand{\frag}[1]{{\itshape «#1»}\par\vspace{1pt}}
% --- Libro 3 · Estrategias ---
\masejemplosde{pu-el-1}{
\ejx{Predecir desde el título}{Un texto se titula «El día que el río cambió de dueño» y su último párrafo empieza: «Desde entonces, la comunidad administra el acueducto». ¿De qué tratará?}{De cómo una comunidad pasó a manejar su propia agua. El título anuncia un cambio y el cierre dice cuál fue: con esas dos piezas ya se sabe qué buscar al leer (quién lo manejaba antes, cómo se logró).}
\ejx{Las primeras frases}{Las primeras frases de tres párrafos son: «El café llegó a Colombia en el siglo XVIII», «Durante el siglo XX fue el principal producto de exportación» y «Hoy, los caficultores enfrentan nuevos retos». ¿Qué estructura tendrá el texto?}{Cronológica: pasado, auge y presente. Las preguntas probablemente pedirán ubicar hechos en una de esas tres épocas.}}
\masejemplosde{pu-el-2}{
\ejx{Tema o idea central}{\frag{Dormir poco no solo produce cansancio. Quien duerme mal recuerda menos, se concentra peor y se enferma con más facilidad. Por eso el sueño debería cuidarse tanto como la alimentación.}¿Cuál es el tema y cuál la idea central?}{Tema: el sueño. Idea central: dormir bien es tan importante para la salud como alimentarse bien. La idea central es una oración completa, no una palabra.}
\ejx{La función de una oración}{En el fragmento anterior, ¿qué función cumple la segunda oración?}{Da las razones (los efectos de dormir mal) que sostienen la conclusión de la tercera. La primera plantea, la segunda argumenta, la tercera concluye.}}
\masejemplosde{pu-el-3}{
\ejx{Tachar}{\frag{El parque ofrece senderos, un lago para remar y un mirador. No tiene restaurante, pero sí zonas de pícnic.}Según el texto, el parque tiene todo lo siguiente EXCEPTO: A) mirador; B) lago; C) restaurante; D) zonas de pícnic.}{A, B y D aparecen: se tachan. El texto dice que no tiene restaurante: C.}
\ejx{Negativa disfrazada}{Con el mismo fragmento: «¿Cuál afirmación es falsa?» A) Se puede remar. B) Hay dónde comer lo que uno lleve. C) Venden almuerzos en el parque. D) Hay senderos.}{«Falsa» funciona igual que EXCEPTO. C contradice «no tiene restaurante»; B es verdadera por las zonas de pícnic.}}
% --- Tipos ---
\masejemplosde{pu-tt-1}{
\ejx{¿Qué predomina?}{\frag{El colibrí bate las alas hasta ochenta veces por segundo. Ese esfuerzo consume tanta energía que debe alimentarse cada pocos minutos.}¿Qué tipo de texto es?}{Expositivo: informa y explica con datos, en tono neutro. No cuenta una historia ni defiende una postura.}
\ejx{El mismo tema, otro tipo}{\frag{Deberíamos sembrar flores nativas en los balcones: sin ellas, los colibríes no tienen de qué vivir en la ciudad.}¿Y este?}{Argumentativo: propone algo («deberíamos») y da una razón. El tema es el mismo; cambió la intención.}}
\masejemplosde{pu-tt-2}{
\ejx{Reconocer el género}{\frag{Señor director: leí con preocupación la noticia sobre el cierre del hospital. Como vecino del barrio, le pido que su periódico siga informando sobre el tema.}¿Qué género es?}{Una carta (al director): tiene destinatario explícito, un remitente que se presenta y una petición.}
\ejx{Noticia o crónica}{\frag{A las cinco de la mañana, cuando todavía no ha salido el sol, doña Rosa ya ha vendido cuarenta arepas. Dice que el secreto es no dormirse.}¿Noticia o crónica?}{Crónica: hay escena, personaje, detalles y voz propia. Una noticia iría directo al hecho, sin recrearlo.}}
\masejemplosde{pu-tt-3}{
\ejx{El esqueleto}{\frag{Las basuras tapaban las alcantarillas y el barrio se inundaba cada invierno. Los vecinos organizaron jornadas de limpieza mensuales. Este año no hubo inundaciones.}¿Qué estructura tiene?}{Problema (inundaciones), solución (jornadas) y resultado. No es solo cronológica: lo que organiza el texto es la respuesta a una dificultad.}
\ejx{Hecho u opinión en un testimonio}{\frag{Trabajé veinte años en la mina. Creo que ningún trabajo enseña tanto sobre el miedo.}¿Qué parte es vivencia y qué parte opinión?}{Vivencia: trabajó veinte años en la mina. Opinión: lo que cree sobre el miedo (la introduce «creo que»).}}
% --- Explícita ---
\masejemplosde{pu-ie-1}{
\ejx{Dos cifras cercanas}{\frag{La biblioteca abrió en 1998 con 2\,000 libros; hoy tiene 20\,000 y recibe unas 300 visitas diarias.}¿Con cuántos libros abrió?}{Con 2\,000. El distractor típico sería 20\,000 (la cifra de hoy) o 1998 (el año).}
\ejx{El orden}{\frag{Antes de sembrar, se abona la tierra; después de sembrar, se riega durante una semana.}¿Qué se hace primero?}{Abonar. «Antes de sembrar» ubica el abono primero, aunque la palabra «sembrar» aparezca antes en la frase.}}
\masejemplosde{pu-ie-2}{
\ejx{Conservar el grado}{\frag{Es probable que el eclipse se vea en buena parte del país.}¿Cuál es la paráfrasis correcta: «el eclipse se verá en todo el país» o «probablemente el eclipse se verá en gran parte del país»?}{La segunda: conserva la probabilidad («probablemente») y la extensión («gran parte»). La primera cambia la probabilidad por certeza y la parte por el todo.}
\ejx{Conservar la relación}{\frag{Las ventas bajaron porque subieron los precios.}¿Equivale a «los precios subieron porque bajaron las ventas»?}{No: invierte causa y efecto, aunque use las mismas palabras.}}
\masejemplosde{pu-ie-3}{
\ejx{Las tres cajas}{\frag{En la encuesta, la mayoría de los jóvenes dijo preferir los libros impresos.}Clasifica cada afirmación en «lo dice», «lo niega» o «no lo dice»: a) más de la mitad de los jóvenes encuestados dijo preferir el papel; b) ningún joven encuestado prefiere los libros impresos; c) los jóvenes leen más que los adultos.}{a) Lo dice: «la mayoría» es más de la mitad. b) Lo niega: la mayoría sí los prefiere. c) No lo dice: el texto no habla de los adultos.}
\ejx{Asociación no es causa}{\frag{Los estudiantes que desayunan suelen tener mejores notas.}¿Dice el texto que desayunar mejora las notas?}{No: dice que las dos cosas suelen ir juntas. La causa podría ser otra (por ejemplo, rutinas más ordenadas en la casa).}}
\masejemplosde{pu-ie-4}{
\ejx{¿Quién opina?}{\frag{Según la alcaldesa, la obra terminará en junio. Los vecinos lo dudan: «Llevamos tres fechas incumplidas», dice uno de ellos.}¿Quién cree que la obra terminará en junio?}{La alcaldesa. El periodista solo lo reporta, y los vecinos piensan lo contrario.}
\ejx{La voz del autor}{En el mismo fragmento, ¿qué afirma el periodista por su cuenta?}{Solo que los vecinos lo dudan. Todo lo demás está atribuido («según», comillas).}}
% --- Vocabulario ---
\masejemplosde{pu-vo-1}{
\ejx{El sinónimo que encaja}{\frag{El equipo logró sortear las dificultades del viaje.}¿«Sortear» equivale a «rifar», «evitar» o «repartir»?}{«Evitar» (esquivar, superar). «Rifar» es otro sentido de la palabra que aquí no cabe: las dificultades no se rifan.}
\ejx{Deducir por el contexto}{\frag{Era un hombre parco: respondía con monosílabos y nunca contaba nada de su vida.}¿Qué significa «parco»?}{De pocas palabras. El propio texto lo explica después de los dos puntos.}}
\masejemplosde{pu-vo-2}{
\ejx{Literal o figurado}{\frag{La noticia cayó como un balde de agua fría.}¿Hubo agua?}{No: es sentido figurado. Significa que la noticia fue una sorpresa desagradable.}
\ejx{El mismo sentido}{«Planta» en «la planta de tratamiento» tiene el mismo sentido que en: a) «regó la planta»; b) «la planta de energía»; c) «la planta del pie».}{b): en ambos casos es una instalación industrial.}}
\masejemplosde{pu-vo-3}{
\ejx{Por la situación}{\frag{Le pedí ayuda con la mudanza y me dijo que sí, pero a la hora de la verdad se hizo el de la vista gorda.}¿Qué hizo?}{Fingió no darse cuenta para no ayudar. La situación (prometió y no cumplió) da el sentido.}
\ejx{Un refrán}{¿Qué refrán resume mejor: «Prometió mucho en campaña y, ya elegido, no hizo nada»?}{«Del dicho al hecho hay mucho trecho»: una cosa es decir y otra hacer.}}
\masejemplosde{pu-vo-4}{
\ejx{Desarmar}{¿Qué significa «hipertensión» si \emph{hiper-} es «exceso»? ¿Y «hipotensión»?}{Tensión (presión) más alta de lo normal; con \emph{hipo-} («por debajo»), más baja de lo normal.}
\ejx{Una palabra nueva}{\frag{El río quedó innavegable después del derrumbe.}¿Qué significa «innavegable»?}{\emph{in-} (no) + navegar + \emph{-able} (que se puede): que no se puede navegar.}}
% --- Referencia ---
\masejemplosde{pu-rf-1}{
\ejx{Género y número}{\frag{Los profesores entregaron los boletines a las familias; estas los recibieron con sorpresa.}¿A quién se refiere «estas»? ¿Y «los»?}{«Estas»: las familias (femenino plural; «profesores» y «boletines» son masculinos). «Los»: los boletines (masculino plural, lo que se recibe).}
\ejx{Un neutro}{\frag{El puente se cayó dos veces en un año. Eso obligó a rediseñarlo.}¿A qué se refiere «Eso»?}{A todo el hecho: que el puente se cayera dos veces. Los neutros recogen ideas completas.}}
\masejemplosde{pu-rf-2}{
\ejx{Reponer el verbo}{\frag{Marta pidió café; su hermano, chocolate.}¿Qué se omite?}{El verbo «pidió»: «su hermano pidió chocolate». La coma marca el hueco.}
\ejx{Una sustitución}{\frag{El jaguar recorre grandes distancias. El felino más grande de América necesita mucho territorio.}¿De quién se habla en la segunda oración?}{Del jaguar: «el felino más grande de América» lo reemplaza para no repetir.}}
% --- Conectores ---
\masejemplosde{pu-cn-1}{
\ejx{Elegir el conector}{«Estudió toda la semana; \rule{12mm}{.4pt}, llegó tranquilo al examen.» ¿«Sin embargo», «por eso» o «aunque»?}{«Por eso»: la segunda idea es consecuencia de la primera.}
\ejx{Lo que cambia un conector}{Compara: «Es joven, pero tiene experiencia» y «Es joven, por eso tiene experiencia».}{La primera tiene sentido: la experiencia contrasta con lo esperado. La segunda no: la juventud no explica la experiencia. El conector revela la lógica.}}
\masejemplosde{pu-cn-2}{
\ejx{Para qué sirve un paréntesis}{\frag{El frailejón (una planta de páramo) retiene el agua.}¿Qué función cumple el paréntesis?}{Aclara qué es el frailejón para el lector que no lo sabe. Si se quita, la oración sigue completa.}
\ejx{Las comillas}{\frag{El «arreglo» de la vía duró dos semanas.}¿Por qué va «arreglo» entre comillas?}{Para marcar ironía: el autor duda de que haya sido un verdadero arreglo.}}
% --- Libro 4 · Inferencia ---
\masejemplosde{pu-in-1}{
\ejx{Con pruebas}{\frag{Julián miró el reloj por tercera vez, guardó el celular y empezó a caminar de un lado a otro del andén.}¿Qué se infiere?}{Que está impaciente o nervioso esperando algo o a alguien. Pruebas: mira el reloj repetidamente y camina de un lado a otro.}
\ejx{Sin pruebas}{Con el mismo fragmento, ¿se infiere que su novia lo dejó plantado?}{No: nada en el texto habla de una novia. Es una historia posible, no una inferencia.}}
\masejemplosde{pu-in-2}{
\ejx{«También»}{\frag{Este año Lucía también ganó la beca.}¿Qué presupone?}{Que alguien más ganó la beca, o que Lucía ya la había ganado antes. «También» siempre suma a algo que se da por sabido.}
\ejx{«Todavía»}{\frag{El ascensor todavía no funciona.}¿Qué presupone?}{Que desde antes no funcionaba y que se espera que funcione.}}
\masejemplosde{pu-in-3}{
\ejx{Directa u original}{\frag{El derrumbe cerró la vía; por eso los camiones no llegaron y el pueblo se quedó sin gasolina.}¿Cuál es la causa directa de la falta de gasolina?}{Que los camiones no llegaron. El derrumbe es la causa original.}
\ejx{Quitar una ficha}{Según el fragmento, ¿qué habría pasado sin el derrumbe?}{Los camiones habrían llegado y habría gasolina: sin la primera ficha, la cadena no cae.}}
\masejemplosde{pu-in-4}{
\ejx{El gesto}{\frag{Cuando anunciaron al ganador, Elena aplaudió despacio, con la mirada fija en el piso.}¿Qué actitud muestra?}{Decepción contenida: aplaude por cortesía, pero el gesto (despacio, mirando al piso) muestra que esperaba otra cosa.}
\ejx{La intensidad}{¿Sería correcto decir que Elena está «furiosa»?}{No: nada indica furia. La opción correcta nombra la emoción con la intensidad que el texto permite.}}
% --- Aplicación ---
\masejemplosde{pu-ap-1}{
\ejx{Llevar el principio}{Un texto afirma: «Lo que se mide, mejora». ¿Cuál caso lo ilustra: a) un corredor que cronometra sus entrenamientos y baja sus tiempos; b) un corredor que entrena con música?}{a): tiene las dos partes del principio (medir y mejorar). b) habla del mismo tema, pero no del mecanismo.}
\ejx{Un contraejemplo}{¿Qué caso contradiría ese principio?}{Alguien que lleva años pesándose a diario sin cambiar ningún hábito y sin mejorar: mide y no mejora.}}
\masejemplosde{pu-ap-2}{
\ejx{Nombrar la relación}{RAÍZ es a ÁRBOL como CIMIENTO es a…}{Edificio: la parte oculta que sostiene al todo.}
\ejx{La relación precisa}{CAPÍTULO es a LIBRO como ESCENA es a: a) actor; b) obra de teatro; c) telón.}{b): parte de un todo que se divide en esas unidades. «Actor» y «telón» pertenecen al teatro, pero no guardan esa relación.}}
% --- Fuentes ---
\masejemplosde{pu-fu-1}{
\ejx{¿Para qué la cita?}{\frag{Suele decirse que los jóvenes no leen. Las cifras de las bibliotecas públicas muestran otra cosa.}¿Para qué menciona el autor lo que «suele decirse»?}{Para oponerse: presenta la idea común y enseguida la contradice con datos.}
\ejx{Apoyo}{\frag{Como advierte la Organización Mundial de la Salud, el sedentarismo es un riesgo creciente.}¿Y aquí?}{Para apoyarse: «como advierte» indica que el autor comparte lo que dice la fuente.}}
\masejemplosde{pu-fu-2}{
\ejx{Pesar dos fuentes}{Sobre la seguridad de un puente opinan el contratista que lo construyó y una universidad que lo revisó sin cobrar. ¿Cuál fuente pesa más?}{La universidad: sabe del tema y no gana nada con el resultado. El contratista tiene un interés directo.}
\ejx{El testimonio}{Un vecino dice: «A mí esa dieta me funcionó». ¿Qué prueba?}{Que a él le funcionó. Una experiencia personal ilustra, pero no permite generalizar.}}
% --- Argumentación ---
\masejemplosde{pu-ag-1}{
\ejx{Separar tesis y razones}{\frag{El recreo no debería acortarse. Los niños que juegan se concentran mejor después, y el juego les enseña a resolver conflictos.}¿Cuál es la tesis y cuáles los argumentos?}{Tesis: el recreo no debería acortarse. Argumentos: mejora la concentración y enseña a resolver conflictos.}
\ejx{Debilitar}{¿Qué dato debilitaría el primer argumento?}{Un estudio que mostrara que los niños rinden igual con recreos más cortos: ataca la premisa de la concentración.}}
\masejemplosde{pu-ag-2}{
\ejx{El tono}{\frag{Otro año, otro informe, otra promesa de que ahora sí se arreglará la vía.}¿Cuál es el tono?}{Escéptico e irónico: la repetición («otro… otra…») y el «ahora sí» muestran que el autor no lo cree.}
\ejx{La intención}{¿Cuál es la intención del fragmento?}{Criticar el incumplimiento repetido, no informar sobre el informe.}}
\masejemplosde{pu-ag-3}{
\ejx{Separar}{\frag{El nuevo parque, que costó 800 millones, es un desperdicio.}¿Qué es hecho y qué es opinión?}{Hecho: costó 800 millones (se puede verificar). Opinión: es un desperdicio.}
\ejx{Un hecho falso}{«Colombia tiene diez océanos.» ¿Hecho u opinión?}{Es del tipo «hecho» (se puede verificar), aunque sea falso. La clasificación no depende de que sea verdad.}}
% --- Figurado ---
\masejemplosde{pu-lf-1}{
\ejx{Reconocer}{\frag{El tiempo es un río que no se devuelve.}¿Qué figura es y qué significa?}{Metáfora (no lleva «como»): el tiempo avanza en una sola dirección y no se recupera.}
\ejx{Dos figuras}{\frag{La tarde bostezaba sobre los tejados, lenta como una tortuga.}¿Qué figuras hay?}{Personificación (la tarde bosteza) y comparación («como una tortuga»). Juntas transmiten lentitud y sopor.}}
\masejemplosde{pu-lf-2}{
\ejx{Detectar el choque}{\frag{Llovió toda la semana de vacaciones. Qué suerte la nuestra.}¿Es irónico?}{Sí: «qué suerte» choca con una semana de lluvia. Quiere decir lo contrario: mala suerte.}
\ejx{Sin ironía}{\frag{Ganamos el partido en el último minuto. Qué suerte la nuestra.}¿Y aquí?}{No: lo dicho coincide con la situación. La misma frase es irónica o no según el contexto.}}
% --- Dos textos ---
\masejemplosde{pu-dt-1}{
\ejx{El punto común}{Texto A: «Los zoológicos son cárceles». Texto B: «Los zoológicos modernos salvan especies». ¿En qué podrían coincidir?}{En que a los dos les importa la suerte de los animales: A la mira en cada animal encerrado y B en la supervivencia de las especies. Tesis opuestas pueden compartir un valor.}
\ejx{La diferencia}{¿En qué difieren exactamente?}{En el juicio sobre los zoológicos: para A son un daño; para B, una herramienta de conservación (y B habla solo de los «modernos»).}}
\masejemplosde{pu-dt-2}{
\ejx{Responder como el otro}{¿Qué respondería el autor del texto B a la frase «son cárceles»?}{Que eso describe a los zoológicos antiguos, no a los modernos que crían especies en peligro: responde desde su tesis.}
\ejx{Una propuesta común}{¿Qué propuesta podrían aceptar los dos?}{Zoológicos con espacios amplios y dedicados solo a conservación: A acepta que mejora el bienestar; B, que mantiene el propósito.}}
